% ─── Chapter 8: Results & Analysis ─────────────────────────────────
\chapter{Results \& Analysis}
\label{ch:results}

This chapter synthesises the experimental outcomes reported in
Chapter~\ref{ch:experiments} into an overall assessment of training
performance.  Results are reported for two distinct settings: \textbf{fixed
teams} (Milestone~1) and \textbf{random teams} (Milestone~2).

\section{Benchmark Metric}

All validation results use the \textbf{benchmark protocol} (3~opponents
$\times$ 50~episodes $=$ 150~total).  The primary metric is the
\textbf{skill score}, a weighted average of per-opponent win rates:

\[
  \text{skill\_score} = \frac{
    1.0 \times \text{WR}_{\text{random}} +
    1.5 \times \text{WR}_{\text{random\_no\_switch}} +
    2.0 \times \text{WR}_{\text{heuristic}}
  }{1.0 + 1.5 + 2.0}
\]

Skill score ranges from 0 (loses all) to 1 (wins all).  The weighting
reflects opponent difficulty: beating the heuristic opponent contributes
twice as much as beating the random opponent.  Per-opponent win rates are
reported with Wilson 95\,\% confidence intervals.

\section{Fixed Teams (Milestone~1)}

With a fixed player team and no gimmicks, the agent achieves strong
performance:

\begin{center}
\small
\begin{tabularx}{0.8\textwidth}{l c}
  \toprule
  \textbf{Metric} & \textbf{Value} \\
  \midrule
  Benchmark skill score        & 0.9911 \\
  Win rate vs.\ random         & 100\,\% \\
  Win rate vs.\ \texttt{random\_no\_switch} & 96\,\%+ \\
  Win rate vs.\ heuristic      & 98\,\% \\
  \bottomrule
\end{tabularx}
\end{center}

% TODO: Insert MLflow benchmark results screenshot for the best fixed-team
% run.  Show per-opponent win rates with Wilson CI bars and the skill score.

This result demonstrates that the agent can learn a highly effective battle
strategy when the team is known and consistent.  The near-perfect performance
against all three opponent tiers indicates that the architecture, observation
pipeline, and reward design are sufficient for fixed-team play.

\subsection{Value Function Improvement}

The reward scaling fix (Section~\ref{sec:reward_scaling}) was the turning
point for the fixed-team setting:

\begin{center}
\small
\begin{tabularx}{\textwidth}{l c c}
  \toprule
  \textbf{Metric} & \textbf{Before Scaling} & \textbf{After Scaling} \\
  \midrule
  Explained variance & ${\sim}0.1$ (sometimes negative) & 0.53--0.87 (median 0.72) \\
  Value loss         & High, no trend                     & 0.08--0.13 (converged by 15K) \\
  Win rate vs.\ heuristic & ${\sim}5$\,\%               & 98\,\% \\
  \bottomrule
\end{tabularx}
\end{center}

% TODO: Insert MLflow comparison chart: before vs after reward scaling(reward scaling 1 instead of 0.1 or 0.05).
% Show explained variance, value loss, and win rate curves side by side.

\subsection{Training Curve}

% TODO: Insert MLflow training curves for the best fixed-team run showing:
% - Win rate per opponent type over time (should show rapid convergence)
% - Value loss convergence (should drop and stabilise early)
% - Entropy decrease as policy becomes confident
% Mark curriculum stage transitions if applicable.

\section{Random Teams (Milestone~2)}

With randomly generated teams for both sides (the standard
\texttt{gen8randombattle} format), the problem is substantially harder.  Each
battle involves unseen teams on both sides(agent can see his own team), requiring the agent to generalise
rather than memorise a single team's strategy. 

\subsection{Current Best Results}

\begin{center}
\small
\begin{tabularx}{0.8\textwidth}{l c}
  \toprule
  \textbf{Metric} & \textbf{Value} \\
  \midrule
  Benchmark skill score (best trial) & 0.3733 \\
  Win rate vs.\ random              & --- \\
  Win rate vs.\ \texttt{random\_no\_switch} & --- \\
  Win rate vs.\ heuristic           & 8\,\% \\
  \bottomrule
\end{tabularx}
\end{center}

% TODO: Fill in the per-opponent win rates for the best random-teams trial
% once the hp sweep completes.  Replace the --- entries above.

The skill score of 0.3733 reflects the difficulty of generalisation: the
agent can beat random opponents reliably but struggles against the heuristic.
This is expected given the high entropy in team composition.  The
hyperparameter sweep (Phase~15) is currently running to optimise
performance in this setting.

\subsection{Hyperparameter Sweep Status}

The Optuna sweep is in progress with the following setup:
\begin{itemize}
  \item 50~trials planned, 600K steps per trial.
  \item Single fixed curriculum stage (no promotion), ensuring comparable
        evaluation across trials.
  \item Objective: benchmark skill score.
\end{itemize}

% TODO: Insert Optuna visualisation once sweep completes:
% - Parallel coordinate plot or importance plot showing which hyperparameters
%   most affect the skill score.
% - Best trial's training curves compared to default config.

\subsection{Noise in Benchmark Scores}

An important observation is that the benchmark skill score has significant
variance due to the stochastic nature of random-team battles.  Lucky team
matchups can inflate the score above what the peak training win rates would
imply.  With only 50~episodes per opponent tier, the Wilson confidence
intervals are wide, particularly for the heuristic tier where win rates are
low.  Future work should increase episode counts for more stable evaluation.

\section{Self-Play Diagnostics}

After the NaN fix (Phase~12), the self-play infrastructure was validated over
500K steps of training:

\begin{center}
\small
\begin{tabularx}{\textwidth}{l c l}
  \toprule
  \textbf{Metric} & \textbf{Value} & \textbf{Interpretation} \\
  \midrule
  Fallback rate       & 0\,\%         & Self-play fully functional \\
  Mapping fallback    & 0\,\%         & Action resolution correct \\
  Top probability     & 0.21 $\to$ 0.47 & Policy becoming confident \\
  Entropy             & 2.0 $\to$ 1.2  & Converging from uniform \\
  Win rate vs.\ random & 91\,\%      & Strong baseline competency \\
  Win rate vs.\ \texttt{random\_no\_switch} & 31\,\% & Room for improvement \\
  Win rate vs.\ self   & 50\,\%      & Near-parity (expected) \\
  \bottomrule
\end{tabularx}
\end{center}

% TODO: Insert MLflow self-play diagnostics charts:
% - Top probability and entropy over training steps
% - Action histogram showing diversification over time

\section{Key Insight: Reward Scaling as the Catalyst}

The central finding from the experimental programme is that \textbf{reward
scaling was the single most impactful change}:

\begin{keyinsight}
For the entire project history, the value function was non-functional
(explained variance ${\sim}0.1$) despite multiple architectural improvements.
A single change, adding \texttt{reward\_scale}$\,=0.1$ to normalise returns
from $[-15,+15]$ to $[-1.5,+1.5]$, caused explained variance to jump to
0.53--0.87 and unlocked competency against heuristic opponents (fixed teams).  Combined with
fixed teams, this produced a benchmark skill score of 0.9911 with 98\,\% win
rate vs.\ heuristic.

The lesson: \textbf{normalise the return scale before tuning anything else.}
The value function must be able to learn before the policy can improve.
\end{keyinsight}

This redirects the optimisation focus for the random-teams setting from
architecture to training dynamics: the same reward-scaling insight applies,
but the higher entropy of random teams requires additional work on
generalisation. This problem remains unsolved for now.

\section{Honest Assessment}

At the time of writing, the project demonstrates:

\begin{center}
\small
\begin{tabularx}{\textwidth}{l X}
  \toprule
  \textbf{What works} & \textbf{What still needs work} \\
  \midrule
  Fixed teams: near-perfect performance (skill score 0.99) &
    Random teams: 0.37 skill score, 8\,\% vs.\ heuristic \\
  Value function learns after reward scaling (EV 0.72) &
    Benchmark score has high variance at 50 eps/tier \\
  Self-play infrastructure functional (0\,\% fallback) &
    LSTM end-to-end validation still pending \\
  Distributed training stable at 48 environments &
    Gauntlet validation pipeline not yet automated \\
  MLflow tracking produces complete, comparable metrics &
    Complete information random vs random as a test \\
  Curriculum learning provides structured progression &
    Team builder (Milestone~3) not yet started \\
  \bottomrule
\end{tabularx}
\end{center}

The engineering foundation is solid and Milestone~1 (fixed teams) is
essentially solved.  The remaining work focuses on Milestone~2 (random teams)
and the team builder for Milestone~3.
